\noindent
Este informe presenta la preparación, auditoría de calidad y análisis descriptivo de \texttt{data/persona.csv}, microdato sociodemográfico oficial de la Encuesta de Hogares de Bolivia considerado por el proyecto. El archivo de origen contiene 39.497 registros y 275 columnas, con una huella criptográfica SHA-256 verificada e inmutable. El resultado analítico se generó como una versión maestra inmutable (\texttt{persona\_clean\_master.parquet} y \texttt{.csv}), reteniendo íntegramente las 39.497 filas y 275 variables, acompañada de 11 vistas temáticas especializadas por población elegible, con trazabilidad exhaustiva en manifiestos y catálogo JSON local con reemplazo atómico, sin dependencia de motores externos como PostgreSQL.

La metodología aplicó reglas de calidad estrictamente respaldadas: validación y tipado de la clave primaria compuesta \texttt{(folio, nro)} (39.497 tuplas únicas, cero duplicados); taxonomía diferenciada de ausencias para respetar los 5,48 millones de tokens \texttt{NA} correspondientes a saltos legítimos del cuestionario; corrección semántica de frontera física en la regla S-01 (reasignación a \texttt{NA} de 3 celdas con jornadas semanales superiores a 168 horas físicas en \texttt{phrs}, \texttt{shrs} y \texttt{tothrs}); y normalización acotada a minúsculas en la regla S-02 (12.516 celdas en 13 variables de respuesta abierta tipo ``Especifique''). Se descartó la eliminación destructiva de variables por umbral del 80\% (regla L-80), tratándose como alerta informativa.

La suite del sistema analítico incluye una aplicación web en Flask que consume la versión validada por catálogo JSON, presentando resúmenes de calidad, perfiles de los 11 universos temáticos y un mapa interactivo de Bolivia con división oficial de los 9 departamentos. Se documenta la distribución exacta observada de la muestra ($N = 39.497$ personas y $N_{\text{hog}} = 12.718$ hogares), delimitando indicadores de salud, educación, empleo e ingresos para uso descriptivo riguroso.
